\documentclass[tikz,border=6pt]{standalone}
\usepackage{ctex}
\usepackage{amsmath}
\usepackage{pgfplots}
\pgfplotsset{compat=1.18}
\usetikzlibrary{calc,angles,quotes,arrows.meta,positioning,decorations.markings}
\begin{document}
% theta = 55 度，第一象限，便于所有线段都可见
\begin{tikzpicture}[scale=5]
  \def\ang{55}
  \pgfmathsetmacro{\cosv}{cos(\ang)}
  \pgfmathsetmacro{\sinv}{sin(\ang)}
  \pgfmathsetmacro{\tanv}{tan(\ang)}
  \pgfmathsetmacro{\secv}{1/cos(\ang)}
  \pgfmathsetmacro{\cscv}{1/sin(\ang)}

  % 坐标轴
  \draw[->,gray!70,thick] (-0.35,0) -- (1.75,0) node[right,black]{$x$};
  \draw[->,gray!70,thick] (0,-0.35) -- (0,1.75) node[above,black]{$y$};

  % 单位圆
  \draw[blue!60!black,very thick] (0,0) circle (1);

  % 角终边（延长到与切线相交，体现 sec/tan）
  \draw[gray!55,thick] (0,0) -- (\ang:\secv);

  % 点 P
  \coordinate (O) at (0,0);
  \coordinate (P) at (\ang:1);
  \coordinate (A) at (1,0);            % (1,0)
  \coordinate (Px) at (\cosv,0);       % P 在 x 轴投影
  \coordinate (T) at (1,\tanv);        % 切线与终边交点 (sec, tan)
  \coordinate (C) at (0,\cscv);        % 终边与 x=... 顶点：csc 在 y 轴

  % --- sin 线段：P 到 x 轴（竖直，红）---
  \draw[red!80!black,very thick] (Px) -- (P);
  % --- cos 线段：O 到 Px（水平，绿）---
  \draw[green!50!black,very thick] (O) -- (Px);

  % --- tan 线段：在 x=1 切线上，A 到 T（橙）---
  \draw[orange!90!black,very thick] (A) -- (T);
  % 切线（虚线，x=1）
  \draw[orange!50,dashed] (1,-0.2) -- (1,1.7);

  % --- sec 线段：O 到 T（沿终边，紫）---
  % 用稍偏移避免与终边重合：直接加粗终边段 O..T
  \draw[purple!70!black,line width=1.2pt,opacity=0.6] (O) -- (T);

  % 点
  \fill[red!80!black] (P) circle (0.012);
  \node[right,font=\scriptsize] at ([xshift=2pt]P) {$P=(\cos\theta,\sin\theta)$};
  \fill[black] (A) circle (0.01);
  \node[below,font=\tiny] at (A) {$(1,0)$};
  \fill[orange!90!black] (T) circle (0.012);

  % 角弧
  \draw[purple!70,thick] (0.22,0) arc (0:\ang:0.22);
  \node[purple!70,font=\small] at (\ang/2:0.31) {$\theta$};

  % 线段标注
  \node[red!80!black,font=\small,right] at ($(Px)!0.5!(P)$) {$\sin\theta$};
  \node[green!45!black,font=\small,below] at ($(O)!0.5!(Px)$) {$\cos\theta$};
  \node[orange!90!black,font=\small,right] at ($(A)!0.5!(T)$) {$\tan\theta$};
  \node[purple!70!black,font=\small,above left] at ($(O)!0.42!(T)$) {$\sec\theta$};

  % 说明框（半透明，放右上空白）
  \node[draw=blue!50!black,fill=white,fill opacity=0.85,text opacity=1,
        rounded corners,font=\scriptsize,align=left,anchor=north west] at (0.95,1.62)
    {单位圆上的线段几何意义\\[1pt]
     纵坐标 $=\sin\theta$（红）\\
     横坐标 $=\cos\theta$（绿）\\
     切线段 $=\tan\theta$（橙）\\
     终边到切线 $=\sec\theta$（紫）};
\end{tikzpicture}
\end{document}
